%!TEX root = ../template.tex
%% ====================================================================
%% Capitulo 7 -- Discussao e proposta metodologica
%% Primeira versao completa 2026-09-08. Linguagem descritiva, sem rotulos
%% internos de organizacao do trabalho. Proveniencia dos numeros:
%%   estimadores robustos:  outputs/T11/runs/20260830T114345550731Z-363f433da25e
%%   bandas honestas:       outputs/T12/runs/20260830T122854929917Z-459cf49c965c
%%   gate de aceitacao:     outputs/T13/runs/20260830T165208669276Z-b583143cacb9
%%   candidatos robustos:   outputs/T14/runs/20260830T165252925802Z-ad4bf78ea405
%%   fichas por baseline:   outputs/T15/runs/20260830T165318477358Z-3d408fcdae7a
%%   referencia congelada:  outputs/PART_C/C1/runs/20260903T002922246855Z-eba4e32107a0
%%   envelopes fixos:       outputs/PART_C/C6/runs/20260903T003533726541Z-892a18a0e064
%%   injecao de desvios:    outputs/PART_C/C5/runs/20260903T004416308950Z-2c5caf3d5988
%%   verdade conhecida:     outputs/PART_C/C8/runs/20260902T165944891648Z-ae3e8022c6a5
%%                          (stress …T224337596785Z-745ef41d27f7, joint …T181655416848Z-46ca4436f02f)
%% Os lotes da exploracao com referencia congelada que a reinspecao de
%% 2026-09-03 classificou como nao finais (comparacao de esquemas de
%% referencia, janelas de reajuste, injecao sobre estados publicados,
%% agregacao por frequencia, eletricidade) NAO sao citados com numeros;
%% onde entram, entram qualitativamente e assinalados. Depois da
%% republicacao (Task 7 do plano de 2026-09-03), substituir onde marcado
%% com %% PENDENTE-REPUBLICACAO.
%% A matematica de cada instrumento vai para o anexo matematico.
%% ====================================================================

\chapter{Discussão e proposta metodológica}
\label{cha:discussao}
% [REV-2026-09-13][C7-01][P2][ESTRUTURA]
\revisaotese{C7-01}{P2}{ESTRUTURA}
{O capítulo mistura discussão, métodos e novos resultados}
{É o capítulo mais longo e introduz candidatos, bandas, critério,
populações, simulador e resultados antes de os discutir. Levar definições
de estimandos, candidatos e simulação para o Cap. 4/Apêndice G; manter
aqui proposta, resultados essenciais da demonstração e interpretação.
Alternativa, se a simulação for contribuição central: capítulo autónomo
de resultados III e discussão mais curta. Recomenda-se primeiro a
redistribuição sem criar outro capítulo, preservando o arco da tese.}%
% [/REV-2026-09-13]

O Capítulo~\ref{cha:diagnostico} estabeleceu que o método de \emph{baselines}
em produção se reproduz exatamente e não é adequado, de forma uniforme, aos
fins que a ISO~50006 lhe atribui. Este capítulo faz três coisas com esse
resultado. Primeiro, lê-o em conjunto: o que os cinco modos de falha têm em
comum, o que a norma exige e não exige, e porque é que a pergunta certa
sobre uma \emph{baseline} de eficiência não é ``prevê bem?'' mas ``revela o
desvio que existe para revelar?''. Segundo, converte o diagnóstico numa
proposta --- um conjunto de linhas orientadoras para construir, aceitar e
manter \emph{baselines} em processo contínuo --- e demonstra-a na própria
unidade: que candidatos da classe linear estática, robustificada e com
bandas honestas, passam um critério de aceitação declarado antes de
qualquer resultado; e o que se aprende quando, em vez de reajustar o modelo
todos os anos, se congela uma referência e se mede tudo o resto como desvio
face a ela, incluindo com verdade conhecida por simulação. Terceiro,
devolve a proposta ao instrumento do Capítulo~\ref{cha:plataforma} como
especificação de um bloco de saúde das \emph{baselines}, e declara os
limites de tudo o que aqui se afirma. A matemática de cada instrumento
consta do Apêndice~\ref{app:matematica}.

\section{Leitura integrada do diagnóstico}
\label{sec:leitura-integrada}

% [REV-2026-09-13][C7-02][P1][ARGUMENTO]
\revisaotese{C7-02}{P1}{ARGUMENTO}
{A discussão amplia indevidamente o veredicto do diagnóstico}
{"Não cumprir nenhuma das três funções" é mais forte que "não adequado
de forma uniforme". O Cap. 6 reconhece conteúdo preditivo no fuel gás,
não decide normalização só pelo MAE e não mede deteção real. Responder
por vetor/uso com evidência e limites, sem converter inadequação parcial
em falha universal. O problema de resolução elétrica também não deriva
da classe linear nem do refit: separar falha de medição, especificação,
transportabilidade e política de referência.}%
% [/REV-2026-09-13]
\paragraph{Reproduzir não é validar.} O primeiro resultado do
Capítulo~\ref{cha:diagnostico} --- a replicação exata das trinta células
vetor--ano --- é a condição de tudo o resto, mas não diz nada sobre a
adequação do modelo: diz que o que se critica é o que está em produção.
Tudo o que se seguiu mostra que um método pode ser reproduzível, documentado
e aplicado com disciplina há seis anos, e ainda assim não cumprir nenhuma
das três funções que lhe são pedidas --- normalizar o consumo pela carga
(a carga só tem conteúdo preditivo no fuel gás), sustentar comparações entre
períodos (a relação anual não é invariante e o veredicto depende de onde
começa o ano) e alimentar alarmes (a banda cobre entre $38\%$ e $99{,}7\%$
dos dias consoante o vetor e o par de anos).

\paragraph{Cinco modos de falha, uma raiz.} Cada vetor falha à sua maneira
--- o fuel gás por deriva paramétrica que a recalibração absorve; o vapor de
24~bar por vacuidade, uma média disfarçada de modelo; o de 10~bar por um
nível em movimento; o de 3~bar por identificação fraca de uma relação que
provavelmente nunca foi linear nem estática; a energia elétrica por não ser
uma medição diária. A diversidade importa porque condena a receita única: um
remédio para a deriva do fuel gás não serve ao 24~bar, onde não há relação a
estabilizar. Mas os cinco partilham uma raiz de desenho, e é sobre ela que a
proposta deste capítulo se constrói: um modelo linear estático, reestimado
anualmente sobre janelas convencionais, aplicado a um processo cujo
comportamento muda mais devagar do que um ano e mais depressa do que a
recalibração o assinala. O mecanismo é o da Secção~\ref{sec:sintese-modos-falha}:
cada janeiro, o reajuste reabsorve nos parâmetros o desvio acumulado no ano
anterior, e o sistema de alarmes, recalibrado sobre o regime já desviado,
volta a tratar como normal o que devia assinalar.

% [REV-2026-09-13][C7-03][P2][ARGUMENTO]
\revisaotese{C7-03}{P2}{ARGUMENTO}
{Evitar a dicotomia fixa = deteta / adaptativa = esconde}
{Uma referência fixa pode prever bem em regime estável; modelos
adaptativos podem detetar mudanças através de estados, parâmetros ou
inovações acumuladas, conforme o desenho. Não é a "única coisa" que
mede desvio. Reformular como tensão entre adaptação e preservação de
uma referência de comparação, sob as políticas avaliadas. Identificar
qual canal carrega o sinal (resíduo, estado, inovação, CUSUM) e que
mudanças foram demonstradas em simulação.}%
% [/REV-2026-09-13]
\paragraph{Previsão e deteção são perguntas diferentes.} O diagnóstico mediu
sobretudo previsão --- \emph{skill} face a referências triviais, erro fora
da amostra, cobertura de bandas --- porque é o que se pode medir sem
calendário operacional. Mas a função de uma \emph{baseline} de eficiência
não é prever o consumo de amanhã: é servir de referência estável contra a
qual um desvio se torna visível. As duas exigências entram em tensão, e a
Secção~\ref{sec:poder-preditivo} já a expôs: uma média móvel de trinta dias
prevê o consumo de vapor melhor do que a relação anual com a carga
precisamente porque incorpora o passado recente --- e, ao incorporá-lo,
incorpora também qualquer desvio persistente, tratando-o como o novo
normal. Um modelo que segue o processo é um bom preditor e um mau detetor;
uma referência congelada é um mau preditor e a única coisa que mede o
desvio. O que este capítulo propõe não é escolher entre os dois: é
reconhecer que são dois objetos com funções distintas, e que o método em
produção tenta ser ambos com um só modelo reajustado uma vez por ano ---
tarde de mais para prever, cedo de mais para detetar.

% [REV-2026-09-13][F17][P2][FIGURA]
\revisaotese{F17}{P2}{FIGURA}
{Esquema central: referência, estado e monitor}
{Após a tabela dos estimandos, inserir diagrama em três ramos:
E\_t,Q\_t,Z\_t -\textgreater{} referência fixa f\_ref -\textgreater{} resíduo;
resíduo -\textgreater{} modelo temporal -\textgreater{} estado estimado + inovação;
resíduo/estado/inovação -\textgreater{} monitor explicitamente definido -\textgreater{} investigação.
Distinguir o que fica fixo e o que se atualiza e indicar quando os
dados entram. Mostrar que o alerta não nasce apenas da escolha do
normalizador. Complementa F07: aqui componentes, lá cronologia.}%
% [/REV-2026-09-13]
\paragraph{Três formas de avaliar, três estimandos.} O que se segue avalia
\emph{baselines} de três maneiras que respondem a perguntas diferentes, e
que convém distinguir antes de ler qualquer número. A primeira, herdada do
Capítulo~\ref{cha:diagnostico}, reajusta o modelo em cada ano e aplica-o ao
seguinte: mede a transportabilidade da relação anual. A segunda, usada na
demonstração das linhas orientadoras (Secção~\ref{sec:demonstracao}),
mantém esses mesmos cinco pares anuais mas substitui o estimador, a banda ou
ambos, e adjudica cada candidato por um critério declarado a priori: mede se
existe um candidato aceitável dentro de uma classe. A terceira
(Secção~\ref{sec:referencia-congelada}) congela uma referência em 2020 e
aplica-a a 2021--2025 sem lhe tocar: mede o desvio face a um ponto fixo, e o
que cada componente de um sistema de monitorização --- normalizador,
escala, monitor, banda --- faz a esse desvio. Os três estimandos são
diferentes por construção, e o mesmo vetor pode ter leituras que só parecem
contraditórias: o vapor de 24~bar é ``estável e calibrado'' na primeira
(declives iguais, banda de produção compatível em cinco pares) e, na
terceira, ``desloca-se de nível sem que nenhum normalizador acrescente nada
à média de 2020''. Não há contradição: a primeira pergunta pela relação com
a carga ano a ano, a terceira pelo nível face a um ano fixo. A
Tabela~\ref{tab:estimandos} fixa os três antes de qualquer resultado.

\begin{table}[htbp]
  \centering
  \caption[Os três estimandos deste capítulo]{As três formas de avaliar uma
  \emph{baseline} usadas na dissertação, o que cada uma mede e o que não
  pode estabelecer.}
  \label{tab:estimandos}
  \small
  \begin{tabular}{p{3.0cm}p{3.6cm}p{3.6cm}p{3.4cm}}
    \toprule
    Avaliação & Desenho & O que mede & O que não estabelece \\
    \midrule
    Diagnóstico (Cap.~\ref{cha:diagnostico}) & modelo do ano $y$ aplicado a $y{+}1$, cinco pares 2020--2025 & transportabilidade da relação anual; calibração da banda; valor de covariáveis & utilidade para deteção; causa da instabilidade \\
    Demonstração (Secção~\ref{sec:demonstracao}) & os mesmos cinco pares; estimador, banda e critério de aceitação declarados a priori & se algum candidato da classe passa o critério & impossibilidade da classe; necessidade de modelos com estado \\
    Referência congelada (Secção~\ref{sec:referencia-congelada}) & referência de 2020 aplicada a 2021--2025; simulação com verdade conhecida & desvio face a um ponto fixo; o que cada componente faz ao desvio & poupança, alarme validado, mecanismo; factos sobre a refinaria a partir da simulação \\
    \bottomrule
  \end{tabular}
\end{table}

\section{Implicações para a conformidade normativa}
\label{sec:implicacoes-normativas}
% [REV-2026-09-13][C7-04][P2][NORMAS]
\revisaotese{C7-04}{P2}{NORMAS}
{Reformular a consequência normativa e discutir literatura}
{Aplicar C2-02/C2-04: não afirmar que a norma certifica só existência,
nem que o pressuposto foi refutado nos cinco vetores (eletricidade foi
excluída e o 24 bar tem resultados não conclusivos para estabilidade).
Apresentar a proposta como operacionalização de adequação/monitorização
com critérios explícitos. Confrontar a contribuição com 3--5 estudos
próximos identificados no Cap. 3; a discussão atual dialoga sobretudo
com os próprios resultados, sem estabelecer novidade relativa.}%
% [/REV-2026-09-13]

A ISO~50001~\cite{ISO50001_2018} obriga a monitorizar, medir e analisar o
desempenho energético; a ISO~50006~\cite{ISO50006_2014} operacionaliza-o em
indicadores e \emph{baselines}, com variáveis relevantes, fatores estáticos
e normalização; a ISO~50015 e o IPMVP~\cite{ISO50015_2014,IPMVP_2022} tratam
da medição e verificação de poupanças. Como o Capítulo~\ref{cha:enquadramento}
mostrou, as três dizem \emph{o que} deve existir e são quase omissas sobre
\emph{como validar} o que existe --- e carregam o pressuposto tácito de que a
relação entre consumo e variáveis relevantes é estável no período de
referência e transportável para o de reporte. O Capítulo~\ref{cha:diagnostico}
é, em rigor, um teste desse pressuposto num processo contínuo, e o resultado
é que ele não se verifica em nenhum dos cinco vetores da unidade.

A consequência normativa é incómoda mas precisa. Um sistema de gestão de
energia pode ser formalmente conforme --- \emph{baselines} definidas,
variáveis relevantes identificadas, indicadores calculados, revisões
anuais documentadas --- e assentar em modelos que não prevêem melhor do que
uma média, cujos parâmetros mudam de sinal entre anos e cujas bandas de
alarme oscilam entre o silêncio e o ruído. A conformidade certifica a
existência do instrumento, não a sua validade; e a norma, ao não exigir
validação, não protege a organização dessa diferença. A ISO~50006 admite
explicitamente a recalibração da \emph{baseline} quando as condições mudam
--- e é precisamente essa admissão que, aplicada como rotina anual, produz o
mecanismo da Secção~\ref{sec:leitura-integrada}: a recalibração sancionada
pela norma é o que esconde o desvio que a norma quer ver.

Daqui decorre a posição desta dissertação: os requisitos de validação têm
de ser acrescentados pela organização, porque a norma não os traz. O
protocolo de avaliação do Capítulo~\ref{cha:metodos} e as linhas
orientadoras que se seguem são uma proposta concreta desses requisitos ---
compatível com o enquadramento normativo, porque não pede nada que a norma
proíba, e mais exigente do que ele, porque pede que uma \emph{baseline} prove
antes de ser usada aquilo que a norma assume.

\section{Linhas orientadoras para a construção e aceitação de \emph{baselines}}
\label{sec:guidelines}

As linhas orientadoras foram formuladas antes da demonstração, revistas em
quatro rondas de revisão adversarial independente (Apêndice~\ref{app:adversarial})
e fixadas na versão que aqui se descreve. Cada uma responde a um modo de
falha do Capítulo~\ref{cha:diagnostico} ou a um defeito apontado na revisão;
nenhuma foi ajustada depois de conhecidos os resultados da
Secção~\ref{sec:demonstracao}. Apresentam-se por ordem de aplicação.

\paragraph{Dados reconciliados como pré-condição.} Nenhuma \emph{baseline} se
estima sobre dados que não passaram a validação de fontes do
Capítulo~\ref{cha:plataforma}: imputações sinalizadas e contadas,
proveniência com \emph{hashes}, e um perfil dos dados --- cobertura,
mecanismo de falta, efeito de qualquer filtro de seleção --- lido e aceite
por uma pessoa antes de qualquer análise. O erro de seleção que invalidou a
análise preliminar desta dissertação (um limiar de carga escrito à mão em vez
de lido da configuração oficial) é o exemplo do que este passo existe para
apanhar.

% [REV-2026-09-13][C7-05][P2][CRITERIOS]
\revisaotese{C7-05}{P2}{CRITERIOS}
{Não converter escolhas prudentes em regras universais}
{Justificar 48--72 h por dinâmica da unidade/eventos; não é um limiar
validado pelos dados disponíveis. Regime quase estacionário e domínio
da regressão são condições diferentes. A exclusão de resposta por erro
instrumental comprovado pode ser legítima: a regra deve distinguir
controlo de qualidade de seleção por desempenho. Conciliar com o filtro
de zeros herdado. Suspender o alarme deste modelo fora do domínio deve
gerar estado "não avaliado", sem sugerir operação normal.}%
% [/REV-2026-09-13]
\paragraph{Domínio definido nas variáveis explicativas, nunca no consumo.} A
\emph{baseline} declara o intervalo de carga em que foi estimada, justificado
por percentis de operação ou limites físicos, e um critério de regime
estacionário (variação de carga abaixo de um limiar; exclusão das primeiras
48 a 72~horas após arranque). A seleção faz-se sobre os regressores
pré-especificados, nunca sobre o consumo --- selecionar sobre a variável
resposta trunca a distribuição do erro e enviesa o modelo~\cite{Heckman1979}. O que a
Secção~\ref{sec:perfil-dados} mostrou --- o limiar oficial remove outro
regime operatório, não uma cauda --- torna esta declaração obrigatória: o
domínio de validade existe de qualquer forma; a diferença está em declará-lo
ou descobri-lo. Na aplicação, um dia fora do domínio é sinalizado como tal
e o seu alarme é suprimido: não é uma anomalia, é uma pergunta a que o
modelo não sabe responder.

\paragraph{Exclusões por causa, nunca por valor.} Um dia é excluído do treino
por um código de razão pré-registado e ligado ao registo de operações
(paragem, arranque, intervenção, campanha), com a declaração de que usos
essa exclusão afeta --- treino, avaliação de desempenho e avaliação de
alarmes são decisões distintas. Nunca por ser um valor extremo. Fica
declarado o risco residual: um registo operacional pode ele próprio ter sido
motivado por um comportamento energético anómalo, e ``documentado'' não é
sinónimo de ``exógeno''.

\paragraph{Especificação antes de robustificação.} Entre o modelo univariado
na carga e qualquer modelo não linear existe o degrau que a ISO~50006 e o
IPMVP sancionam~\cite{ISO50006_2014,IPMVP_2022}: a regressão linear múltipla com variáveis relevantes
pré-especificadas (densidade da carga, condições de \emph{flash}, sinais de
modo de operação). A Secção~\ref{sec:suficiencia} mostrou o que este degrau
vale na unidade --- consistente num vetor, prejudicial noutros --- e daí a
regra: uma covariável é admitida só por valor preditivo incremental
demonstrado fora da amostra, para a frente no tempo, com um critério que
inclua a banda e não apenas o erro pontual; nunca por catálogo.

\paragraph{Estimação robusta, pré-especificada.} O estimador primário é
declarado antes de ver os dados --- um estimador-M de Huber~\cite{Huber1964,HuberRonchetti2009} com constante
de afinação fixa ($c = 1{,}345$~\cite{HollandWelsch1977}) e escala pela
mediana dos desvios absolutos ---
com os mínimos quadrados reportados ao lado, e um limiar de divergência
declarado entre os dois que, se excedido, obriga a investigar: a divergência
pode significar má especificação, não valores atípicos. A alavancagem
trata-se separadamente, porque um estimador de Huber não resiste a pontos de
alavancagem elevada~\cite{RousseeuwLeroy1987}. O estimador de Theil--Sen~\cite{Sen1968} serve de referência livre de
distribuição em regressão simples. Um ponto e uma banda robustos são
problemas diferentes, e a linha seguinte trata do segundo.

\paragraph{Bandas honestas, com três desfechos possíveis.} A banda de alarme
constante $\pm 2\hat\sigma$ assume resíduos independentes e homocedásticos que
a Secção~\ref{sec:diagnosticos-desc} mostrou não existirem. A rota primária
modela explicitamente a dependência residual --- um processo autorregressivo
sobre os resíduos~\cite{BoxJenkins2015}, com desfasamentos civis exatos e
ordem escolhida só no treino --- e produz intervalos condicionais, a um passo, validados pela
cobertura empírica no ano seguinte. A rota de sensibilidade é a regressão
quantílica~\cite{KoenkerBassett1978}, com restrições de não
cruzamento~\cite{Chernozhukov2010} e incerteza por \emph{bootstrap}
de blocos --- sobre uma janela plurianual, nunca sobre um ano ($n \approx
300$ dias a $\tau = 0{,}025$ deixa cerca de sete observações em cada cauda).
O ponto essencial é o terceiro desfecho: a avaliação pode concluir que
\emph{nenhuma} rota é suportada para um vetor, e esse veredicto é
reportável e bloqueia o alarme por banda nesse vetor. Uma banda não
validada não é uma banda pior; é a ausência de banda.

% [REV-2026-09-13][C7-06][P1][ACEITACAO]
\revisaotese{C7-06}{P1}{ACEITACAO}
{Compatibilidade por não-rejeição não valida uma banda}
{O critério aceita cobertura quando o IC contém a nominal: um IC muito
largo facilita passar, premiando falta de precisão. Não chamar a isso
validação/calibração demonstrada. Preservar a adjudicação pré-especificada
como resultado histórico e designar "compatível pelo critério adotado".
Para uma versão futura, fixar margens de cobertura/materialidade com
a operação, precisão mínima e desfecho inconclusivo; avaliar equivalência
ou não-inferioridade apropriada, largura e desempenho de deteção.
Não escolher margens depois de ver quais fazem passar candidatos.
Fonte: \url{https://doi.org/10.1177/1948550617697177}}%
% [/REV-2026-09-13]
% [REV-2026-09-13][C7-07][P1][ACEITACAO]
\revisaotese{C7-07}{P1}{ACEITACAO}
{A regra 4/5 precisa de justificação e avaliação conjunta}
{Pré-especificação evita ajustar a regra ao resultado, mas não a torna
adequada. Explicar porquê 4/5, custo de aceitar/rejeitar e dependência entre
pares. Os dois membros podem cumprir 4/5 em anos diferentes, deixando
só 3/5 com ambos: declarar se esse é o critério pretendido.
Exigir ganho face à média exclui por construção referências constantes
potencialmente úteis à monitorização. Separar aceitação de normalizador
dependente de carga, preditor e detetor. Simular erros de decisão do
procedimento antes de o recomendar para entrada em produção.}%
% [/REV-2026-09-13]
\paragraph{Critério de aceitação com validação temporal encadeada.} Um
candidato --- estimador, banda e domínio --- só entra em produção depois de
passar um critério declarado antes de qualquer resultado, avaliado para a
frente no tempo. A formulação original exigia uma partição selada: construção
e afinação em 2020--2023, seleção congelada em 2024, uma única avaliação em
2025 intocado. Foi substituída, por decisão datada, por validação temporal
encadeada --- os mesmos cinco pares anuais do diagnóstico ---, porque 2025 e
o parcial de 2026 já tinham informado o Capítulo~\ref{cha:diagnostico} e
nenhum resultado obtido depois disso pode chamar-se \emph{holdout}. O critério
tem dois membros e aceita um candidato apenas se ambos se cumprem: conteúdo
preditivo --- ganho sobre a média do ano de treino, com o intervalo de
confiança a excluir zero, em pelo menos quatro dos cinco pares (é o que
impede que um modelo vazio, com $R^2 \approx 0$, seja aceite por ter uma
banda razoável, o defeito da primeira versão das linhas orientadoras); e
banda suportada --- cobertura compatível com a nominal, pelo intervalo, em
pelo menos quatro dos cinco pares. O limiar de quatro em cinco foi fixado
antes de qualquer resultado e não se mexe depois; a
Secção~\ref{sec:demonstracao} mostra o que depende dele.

\paragraph{Recalibração por evidência, com políticas separadas.} Gerar um
alarme, investigar, suspender temporariamente um modelo, recalibrar e
verificar a recalibração são cinco decisões com políticas distintas e
declaradas. A recalibração exige uma mudança detetada --- um sinal de um
monitor sequencial~\cite{Page1954,Lai1995} sobre as inovações, não uma data
no calendário --- mais
uma revisão documentada; mas a revisão pode concluir ``causa desconhecida,
recalibrar com sinalização'', porque exigir causa demonstrada
incondicionalmente atrasaria a adaptação exatamente quando o sistema muda.
Os mesmos resíduos nunca servem para declarar a mudança e para certificar o
modelo recalibrado: a certificação usa só dados posteriores. É esta linha
que substitui a recalibração anual de rotina, e é a resposta direta ao
mecanismo da Secção~\ref{sec:leitura-integrada}.

\paragraph{Mínimos de dados e governança.} Não se estima uma \emph{baseline}
se o número de observações está abaixo de um mínimo declarado (formalizando
o parâmetro que o método de produção já tem), se o treino cobre menos do que
uma fração declarada do domínio de aplicação, ou se a fração imputada excede
um teto. A governança --- versões, registo de alterações, aprovação nominal,
revisão agendada --- e uma ficha de uma página por \emph{baseline} (domínio,
janela, exclusões com códigos de razão, estimador, método de banda e o seu
desfecho, resultados da aceitação, \emph{hash} de proveniência) fecham o
conjunto. A ISO~50001 é uma norma de sistema de gestão; sem governança, as
linhas orientadoras não são adotáveis.

\section{Demonstração na unidade}
\label{sec:demonstracao}
% [REV-2026-09-13][C7-18][P1][ESTIMANDO]
\revisaotese{C7-18}{P1}{ESTIMANDO}
{O critério de aceitação tem de corresponder ao uso da baseline}
{O gate penaliza cobertura baixa no ano seguinte, mas a referência
congelada interpreta essa mesma queda como sinal a preservar. A tensão
é substantiva: sem eventos, não se distingue mudança real de falha
preditiva. Definir se se aceita um preditor sob transportabilidade, um
normalizador ou um detetor sob hipótese de ausência de mudança.
A aceitação do primeiro não certifica o terceiro, nem a rejeição do
primeiro invalida a utilidade da referência. Explicitar isto antes da
tabela dos 32 candidatos e no desenho das duas pistas da plataforma.}%
% [/REV-2026-09-13]

As linhas orientadoras foram demonstradas nos quatro vetores diários da
unidade, sobre os mesmos cinco pares anuais 2020${\to}$2021 a
2024${\to}$2025 do Capítulo~\ref{cha:diagnostico}, com constantes fixadas a
priori ou afinadas apenas dentro do ano de treino, conclusões só por
exclusão do intervalo de confiança por \emph{bootstrap} de blocos móveis, e
margens materiais declaradas como não elicitadas --- não há, nesta unidade,
um limiar acordado a partir do qual um desvio importa. Nenhuma linha desta
secção se descreve como \emph{holdout}.

\subsection{Estimadores robustos}
\label{sec:demo-robustos}

Para cada vetor e par anual estimaram-se três retas sobre o ano de treino
--- mínimos quadrados (o método de produção), Huber e Theil--Sen --- e
compararam-se as perdas absolutas diárias no ano seguinte, sem reestimação,
por contraste emparelhado $|e_{\mathrm{OLS}}| - |e_{\mathrm{rob}}|$ com
intervalo de confiança a 95\%. São quarenta contrastes.

\begin{table}[htbp]
  \centering
  \caption[Estimadores robustos contra mínimos quadrados]{Contrastes
  emparelhados de perda absoluta no ano seguinte, quatro vetores $\times$
  cinco pares. ``Robusto melhor'' e ``OLS melhor'': o intervalo MBB a 95\%
  exclui zero do lado respetivo.}
  \label{tab:robustos}
  \begin{tabular}{lccc}
    \toprule
    Estimador & Robusto melhor & OLS melhor & Sem evidência \\
    \midrule
    Huber ($c = 1{,}345$) & 11 & 6 & 3 \\
    Theil--Sen            & 7  & 8 & 5 \\
    \bottomrule
  \end{tabular}
\end{table}

Nenhum estimador domina (Tabela~\ref{tab:robustos}). O Huber ganha em onze
dos vinte contrastes e perde em seis; o Theil--Sen é mais extremo, sete
contra oito. O maior ganho é do Theil--Sen no vapor de 3~bar em
2022${\to}$2023 ($+4{,}68$~t~GNE/d, $+15{,}3\%$). Três leituras
qualificam a contagem. Contar vitórias não é agregar: no vapor de 24~bar o
Huber ganha três pares em cinco e tem, ainda assim, o erro absoluto médio
agregado ligeiramente pior. A explicação intuitiva --- ``a robustez ajuda
onde há caudas pesadas'' --- é contradita pelos próprios dados: o vapor de
3~bar, onde o Huber mais ganha, é o vetor com menor curtose residual, e o
fuel gás de 2024, com curtose $15{,}7$, não beneficia. E melhor ponto não é
melhor banda: no vapor de 24~bar em 2024${\to}$2025 o Huber reduz o erro
médio em $2{,}9\%$ e a banda marginal construída sobre a sua escala robusta
cobre $48{,}5\%$ dos dias, contra $93{,}9\%$ da banda de produção --- a
escala robusta é $0{,}55$ da escala dos mínimos quadrados, e uma banda mais
estreita em torno de um ponto marginalmente melhor é uma banda muito pior:
a nitidez só vale sujeita à calibração~\cite{Gneiting2007}.
A robustez na estimação e a robustez do detetor são problemas distintos, e
esta é a razão de a linha orientadora os separar.

\subsection{Bandas honestas}
\label{sec:demo-bandas}
% [REV-2026-09-13][C7-08][P2][MODELOS]
\revisaotese{C7-08}{P2}{MODELOS}
{Nomear a construção e tratar a quantílica anual como sensibilidade}
{Preferir "bandas condicionais e quantílicas" a "honestas", que sugere
garantia ainda não demonstrada. A linha orientadora prescreve janela
plurianual, mas a demonstração executa quantílica anual: identificar a
diferença de protocolo e não contar esse resultado como teste da
recomendação plurianual. Explicar fallback para banda marginal e fração
de dias afetados; comparar larguras/coberturas na mesma população e
com o mesmo nível nominal.}%
% [/REV-2026-09-13]

A rota primária ajustou, sobre os resíduos dos mínimos quadrados de cada ano
de treino, um processo autorregressivo com desfasamentos civis exatos e
ordem escolhida pelo critério de informação bayesiano~\cite{Schwarz1978} só
no treino (entre um e sete dias), com predição a um passo no ano seguinte e recurso declarado à banda
marginal quando o histórico não chega; produziu duas construções, uma
gaussiana ($\pm 2\hat\sigma$ das inovações, primária no veredicto) e uma por
quantis empíricos das inovações. A rota de sensibilidade foi a regressão
quantílica a $\tau = 0{,}02275$ e $0{,}97725$ --- os quantis que
correspondem a $\pm 2\sigma$ numa normal --- com verificação de não
cruzamento que recusa o resultado em vez de o corrigir. Uma banda é
suportada num vetor se o intervalo de confiança da sua cobertura inclui
$95{,}45\%$ em pelo menos quatro dos cinco pares, sem degeneração.

\begin{table}[htbp]
  \centering
  \caption[Bandas alternativas nos vinte pares]{Cobertura no ano seguinte
  das quatro construções de banda, vinte pares vetor--transição. ``Inclui'':
  o intervalo MBB da cobertura inclui a referência de $95{,}45\%$; ``sub'' e
  ``sobre'': exclui-a por baixo ou por cima; NE: não estimável (não
  convergência declarada). Meia-largura mediana relativa à banda de
  produção.}
  \label{tab:bandas}
  \begin{tabular}{lccccc}
    \toprule
    Banda & Inclui & Sub & Sobre & NE & Meia-largura \\
    \midrule
    Produção $\pm 2\hat\sigma$   & 10 & 7  & 3 & 0 & $1{,}00$ \\
    AR, gaussiana                & 9  & 6  & 5 & 0 & $0{,}63$ \\
    AR, quantis empíricos        & 10 & 6  & 4 & 0 & $0{,}59$ \\
    Quantílica anual             & 5  & 12 & 2 & 1 & $0{,}97$ \\
    \bottomrule
  \end{tabular}
\end{table}

O veredicto é que \textbf{nenhuma rota candidata é suportada em nenhum dos
quatro vetores} (Tabela~\ref{tab:bandas}) --- e é preciso dizê-lo com
exatidão, porque a frase ``nenhuma banda é suportada'' seria falsa: a banda
de produção do vapor de 24~bar é compatível com a referência nos cinco pares.
O que a rota autorregressiva faz é instrutivo nas duas direções. Recupera
as falhas mais graves da produção --- no vapor de 10~bar em 2024${\to}$2025
a cobertura passa de $52{,}9\%$ para $95{,}6\%$ (mais $42{,}7$ pontos,
intervalo $[30{,}2;\,56{,}3]$); no de 3~bar em 2021${\to}$2022, de $37{,}8\%$
para $76{,}1\%$ --- e quebra três dos cinco sucessos do 24~bar, porque a sua
banda, $0{,}63$ da largura da produção, é estreita de mais onde a de
produção já era adequada. O saldo nos vinte pares é de menos um. Dois
pormenores mudam a leitura. A rota autorregressiva não é só uma banda: ao
usar os resíduos dos dias anteriores, muda o \emph{centro} da predição ---
é um modelo diferente, com memória, e parte do que recupera é o nível que os
mínimos quadrados anuais não seguem (o que a Secção~\ref{sec:referencia-congelada}
retoma). E a regressão quantílica sobre um ano tem sete ou oito observações
por cauda: o seu resultado negativo é da implementação anual, que a linha
orientadora já desaconselhava, e não da rota.

\subsection{O critério de aceitação e o que decidiu}
\label{sec:demo-gate}
% [REV-2026-09-13][C7-09][P1][ESTIMANDO]
\revisaotese{C7-09}{P1}{ESTIMANDO}
{O candidato com AR muda a previsão; o ganho tem de corresponder}
{O texto adjudica conteúdo preditivo com E2 (reta anual), mas reconhece
que a construção AR muda o centro usando resíduos passados. Se o
candidato é a previsão completa reta+AR, medir também o seu MAE contra
comparadores nos mesmos dias e com a mesma informação disponível.
Se se pretende validar separadamente normalizador e monitor, explicitar
dois objetos e dois critérios. Até clarificar, zero candidatos aceites
refere-se à regra de adjudicação adotada, não à inutilidade da previsão
completa AR. Conferir implementação e não apenas renomear a tabela.}%
% [/REV-2026-09-13]
% [REV-2026-09-13][F18][P2][FIGURA]
\revisaotese{F18}{P2}{FIGURA}
{Matriz de decisão por candidato e por par anual}
{Depois da definição do candidato, apresentar matriz com linhas
estimador+banda+vetor, colunas cinco transições e duas marcas por célula
(ganho/cobertura). Distinguir rejeitado, compatível, inconclusivo e NE.
À direita, decisão 4/5; em painel secundário, sensibilidade 3/5.
Dar nomes às oito combinações efetivamente avaliadas por vetor:
3 estimadores x 6 bandas x 4 vetores não é um fatorial de 32 células.
A tabela atual resume bem; a matriz completa pode ir ao Apêndice D.}%
% [/REV-2026-09-13]

O critério foi aplicado por adjudicação pura sobre os resultados já
publicados --- o poder preditivo do Capítulo~\ref{cha:diagnostico} e as
bandas da secção anterior ---, sem recalcular nada, com cada valor lido de
uma corrida identificada e verificada por \emph{hash}. Um candidato é um par
(estimador, banda); com os mínimos quadrados e as quatro construções de
banda, são dezasseis decisões.

\begin{table}[htbp]
  \centering
  \caption[Critério de aceitação: candidatos de mínimos quadrados]{Adjudicação
  dos dezasseis candidatos com estimador de mínimos quadrados. Conteúdo
  preditivo: pares (em cinco) com ganho sobre a média do treino e intervalo a
  excluir zero; banda: pares com cobertura compatível com a nominal. Aceite
  se ambos $\geq 4/5$.}
  \label{tab:gate}
  \small
  \begin{tabular}{lcccccc}
    \toprule
    Vetor & Cont. preditivo & Produção & AR gauss. & AR emp. & Quantílica & Aceites \\
    \midrule
    Fuel gás     & \textbf{4/5} & 3/5 & 3/5 & 3/5 & 2/5 + NE & 0/4 \\
    Vapor 24 bar & 1/5 & \textbf{5/5} & 2/5 & 2/5 & 2/5 & 0/4 \\
    Vapor 10 bar & 0/5 & 1/5 & 3/5 & 3/5 & 0/5 & 0/4 \\
    Vapor 3 bar  & 2/5 & 1/5 & 1/5 & 2/5 & 1/5 & 0/4 \\
    \bottomrule
  \end{tabular}
\end{table}

\textbf{Nenhum dos dezasseis candidatos é aceite} (Tabela~\ref{tab:gate}), e
as razões separam-se limpamente por vetor. O fuel gás cai só pela banda:
tem conteúdo preditivo em quatro dos cinco pares, mas nenhuma construção o
acompanha com cobertura compatível em mais de três. Os vapores caem pelo
conteúdo preditivo --- um, zero e dois pares em cinco --- e a única célula
em que uma banda passa, a de produção do 24~bar, protege um modelo sem
relação a proteger: é o espelho do fuel gás, uma banda suportada sem modelo
por trás.

Os candidatos robustos completos --- Huber e Theil--Sen, cada um com a banda
marginal sobre a sua escala robusta e com a rota autorregressiva sobre os
seus resíduos, mais dezasseis decisões --- não mudam o desfecho: zero
aceites. Os quatro candidatos do fuel gás mantêm o conteúdo preditivo
(quatro em cinco com ambos os estimadores) e nenhuma banda passa de três em
cinco; os vapores continuam sem conteúdo preditivo (o melhor, Theil--Sen no
3~bar, chega a três em cinco). Trocar o estimador pontual não resgata a
classe. A correção autorregressiva sobe o número de pares compatíveis de
nove, com a banda marginal, para dezoito --- sem nunca chegar a quatro num
mesmo candidato; o caso extremo é o vapor de 24~bar com Theil--Sen em
2024${\to}$2025, em que a banda marginal cobre $14{,}6\%$ e a autorregressiva
$96{,}6\%$.

O que o zero em trinta e dois significa, e o que não significa, tem de ficar
escrito com o mesmo cuidado. Significa que nenhuma das combinações avaliadas
--- três estimadores, seis construções de banda, quatro vetores --- cumpre
o critério declarado sob a regra de quatro em cinco. Não demonstra a
impossibilidade da classe linear, nem a necessidade de modelos com estado
--- essa frase chegou a constar das fichas e foi retirada. E o resultado é
sensível ao limiar: com uma regra de três em cinco seriam três candidatos
aceites em dezasseis; a regra de quatro em cinco estava pré-registada e não
se altera depois de conhecido o resultado, mas o leitor deve saber que a
decisão assenta em cinco pares anuais, e que a potência de uma regra binária
sobre cinco observações dependentes nunca foi calculada (Secção~\ref{sec:limitacoes}).
O que a demonstração estabelece de sólido é o \emph{procedimento}: um
critério declarado a priori, adjudicação separada do cálculo, razões de
rejeição distintas por vetor, e um resultado negativo que fica publicado
como tal em vez de ser ajustado até aparecer um vencedor.

\subsection{Fichas por \emph{baseline}}
\label{sec:demo-fichas}

A última linha orientadora foi demonstrada produzindo, para cada um dos
quatro vetores, a ficha de uma página que a governança exige: domínio de
validade, janela, exclusões, estimador, método de banda e o seu desfecho,
resultados do critério de aceitação, e proveniência --- setenta e seis
campos por vetor, cada um lido de uma corrida identificada e verificado por
\emph{hash} contra o manifesto de origem, sem nenhum valor escrito à mão. As
fichas não calculam nada: consolidam. E é exatamente por isso que servem de
especificação para a plataforma (Secção~\ref{sec:implicacoes-plataforma}): o
que uma pessoa precisa de saber para confiar numa \emph{baseline} cabe numa
página, desde que cada campo tenha origem rastreável. Duas erratas
apanhadas em revisão --- um máximo de 2026 publicado como se fosse do
período 2020--2025, e uma frase sobre ``deriva de nível como falha estrutural
documentada'' que as corridas não sustentam --- foram corrigidas na versão
final e ficam registadas no Apêndice~\ref{app:adversarial}.

\section{Referência congelada e verdade conhecida}
\label{sec:referencia-congelada}

A demonstração anterior manteve a perspetiva do método de produção:
reajustar todos os anos e perguntar se o ajuste transporta. A exploração
que se segue muda de perspetiva. Congela uma \emph{referência} --- a relação
estimada em 2020, com a sua escala de erro e, quando aplicável, o seu estado
inicial --- e aplica-a a 2021--2025 sem lhe tocar. Tudo o que se mede na
aplicação passa a ser desvio face a 2020, e o desenho separa três objetos
que o método de produção mistura num só: o resíduo da referência (o desvio
físico face ao ponto fixo), o estado filtrado (o que um modelo adaptativo
segue) e a inovação (o que o filtro não previu). O reajuste anual é aqui
uma sensibilidade, não o primário. Todos os resumos usam uma população
declarada --- dias de 2021--2025 retidos pela seleção oficial, dentro do
suporte de carga de 2020, com observação e predição finitas --- e uma
comparação entre modelos faz-se sempre na interseção dessas populações; foi
a imposição desta regra, numa revisão independente, que inverteu várias
ordenações da primeira versão desta exploração, e a versão que aqui se cita
é a corrigida.

\subsection{O que muda quando não se reajusta}
\label{sec:referencia-fixa}

% [REV-2026-09-13][C7-10][P2][POPULACAO]
\revisaotese{C7-10}{P2}{POPULACAO}
{A interseção de 876 dias altera a população de interesse}
{Quantificar 876 como fração dos dias candidatos e primários, por ano
e vetor; identificar que modelo/requisito elimina cada grupo.
Comparação na interseção é justa entre modelos, mas pode privilegiar
períodos fáceis e limitar a aplicação operacional. Mostrar resultados
na população de cada modelo como sensibilidade separada, sem misturar
denominadores. Justificar 2020 como referência disponível, não regime
limpo/eficiente comprovado; avaliar outra referência como sensibilidade
depois de corrigir os lotes, nunca escolhendo a de menor consumo.}%
% [/REV-2026-09-13]
\paragraph{A referência congelada mede o desvio.} Aplicada a 2021--2025 na
população comum aos cinco normalizadores avaliados ($n = 876$ dias), a
relação de 2020 tem, no fuel gás, \emph{skill} de $41$ a $54\%$ face à média
de 2020 --- os normalizadores continuam a bater a média cinco anos depois.
No vapor de 24~bar nenhum normalizador acrescenta nada à média de 2020 (de
$-5{,}6\%$ a $+0{,}3\%$); no de 10~bar o ganho é de $2$ a $6\%$; no de
3~bar a média de 2020 é o melhor ``modelo'' de todos, por $16$ a $25\%$. O
desvio padronizado médio face à referência é de $0{,}97\sigma$ no fuel gás
e $1{,}73\sigma$ no vapor de 10~bar: a aplicação já não se parece com o ano
de referência, e um normalizador congelado em 2020 é o que mede essa
distância. Uma banda de $95\%$ calibrada em 2020 cobre, em 2021--2025 e nos
dias em que a comparação é admissível, entre $0{,}564$ e $0{,}984$ dos dias
consoante o vetor, o modelo e a regra --- mediana entre vetores de
$0{,}71$ a $0{,}86$ por modelo ---, com o vapor de 24~bar acima de $0{,}94$
em todos os tercis de carga e o de 10~bar a cair a $0{,}43$ no tercil de
carga baixa. Isto é a medição, não um defeito a corrigir: a banda diz que a
aplicação já não se parece com a referência, e nenhuma banda é alargada até
o número voltar. Perder cobertura nominal face a uma referência congelada
não é, por si, um veredicto de ineficiência (Secção~\ref{sec:limitacoes});
é a quantidade que um sistema de monitorização deve mostrar e o método de
produção reabsorve.

\paragraph{Um controlo adaptativo cobre mais, por seguir o desvio.} Nos
mesmos dias, um controlo por quantis móveis dos noventa dias civis
anteriores --- só com passado admissível --- cobre $0{,}87$ a $0{,}91$ dos
dias nos quatro vetores, contra $0{,}60$ a $0{,}99$ da banda fixa; e fá-lo
com uma banda mais estreita, em unidades da escala de 2020, no vapor de
10~bar ($3{,}02$ contra $4{,}48$) e no de 24~bar, e mais larga no fuel gás.
Cobre mais porque se adapta ao desvio que a banda fixa existe para revelar
--- é a Secção~\ref{sec:leitura-integrada} em números --- e por isso não é
um substituto: é um controlo, pontuável em $38\%$ dos dias primários
(exige sessenta erros admissíveis nos noventa dias anteriores).

% [REV-2026-09-13][F19][P2][FIGURA]
\revisaotese{F19}{P2}{FIGURA}
{Visualizar quanto sinal cada componente preserva}
{Inserir junto da experiência de injeção uma figura com três painéis
no mesmo eixo de dias: sinal conhecido (degrau/rampa); resíduos da
referência fixa e do refit; estado, inovação e estatística do monitor.
Marcar treino, início do desvio, reajuste, lacunas e limiar de deteção.
Curvas do lote validado e mesma população; legenda distingue identidade
algébrica, atenuação medida e deteção. Não usar uma imagem conceptual
para provar os valores 0,136/0,002 ou 3\%.}%
% [/REV-2026-09-13]
% [REV-2026-09-13][C7-11][P1][ALGEBRA]
\revisaotese{C7-11}{P1}{ALGEBRA}
{Reajustar não apaga universalmente o sinal de mudança}
{A frase "reajustar é apagar" excede a experiência. Para OLS com
intercepto, um degrau constante em todo o conjunto usado no refit é
absorvido pelo intercepto; uma mudança que afeta só parte da janela
não desaparece necessariamente, nem há garantia sobre dados futuros.
Especificar fração de preservação: contraste injetado menos controlo,
canal, período e normalização; médias de resíduos podem cancelar sinal.
Reportar atenuação sob este desenho, não impossibilidade geral de deteção.}%
% [/REV-2026-09-13]
\paragraph{Um desvio inserido sobrevive à referência fixa e não sobrevive ao
reajuste.} Para medir o mecanismo diretamente, inseriu-se um desvio
conhecido --- um degrau e uma rampa, em unidades da escala de cada vetor
--- sobre os dados reais de aplicação, e mediu-se a fração desse desvio que
sobrevive em cada canal. O resíduo da referência congelada preserva-o por
inteiro, em qualquer das três leituras (janela inteira, dias em amplitude
plena, primeira quinzena em amplitude plena) --- é uma identidade algébrica,
não um mérito, e é o que ``referência'' significa. Um reajuste da referência
sobre o período desviado apaga um degrau por completo e deixa passar $3\%$
de uma rampa. E a inovação de um filtro de nível~\cite{DurbinKoopman2012} assimila um
desvio sustentado em poucos dias: da primeira quinzena de um degrau retém
$0{,}136$; da janela inteira, $0{,}002$; uma rampa nunca chega a produzir
inovação, porque sobe mais devagar do que o filtro aprende. Esta última
leitura é sobre a inovação de um filtro, não sobre uma estatística de
monitorização --- se um monitor deteta a mudança é a pergunta da secção
seguinte. O que fica estabelecido aqui é o mecanismo do
Capítulo~\ref{cha:diagnostico}, agora medido: reajustar é apagar.

\paragraph{Reajustar com janelas curtas segue melhor onde há o que seguir.}
% [REV-2026-09-13][C7-12][P1][EVIDENCIA]
\revisaotese{C7-12}{P1}{EVIDENCIA}
{Um lote inválido também não sustenta uma conclusão qualitativa}
{O próprio capítulo diz que estes lotes comparam populações diferentes
ou comprimem calendário. Omitir números não resolve o problema:
a ordenação "janela curta melhor" também depende do estimando.
Retirar a conclusão do corpo até republicação/verificação; ou descrevê-la
apenas como hipótese e trabalho pendente, sem a usar para justificar
recomendações. Manter a mesma regra para todas as partes em revisão.}%
% [/REV-2026-09-13]
%% PENDENTE-REPUBLICACAO: numeros das janelas de reajuste (lote C5-refit)
%% so depois da republicacao da Task 7; ate la, leitura qualitativa.
A sensibilidade que reajusta a referência com janelas móveis de seis, doze e
vinte e quatro meses, ou expansiva~\cite{PesaranTimmermann2007}, dá uma
leitura qualitativa que não
depende dos números finais (em revisão à data desta versão): no vapor de
10~bar --- o vetor com o desvio de $1{,}73\sigma$ --- a janela mais curta
tem o menor erro e o maior salto de parâmetros entre reajustes, e o erro
cresce com o comprimento da janela; no fuel gás a ordem inverte-se, e a
janela expansiva tem o menor erro e o menor salto, porque não há desvio que
compense o ruído de reajustar com poucos dados. Quem segue melhor é quem
salta mais, mas só onde há o que seguir --- e um sistema que escolhesse a
janela pelo erro escolheria, no 10~bar, exatamente o modelo que mais depressa
esconde o desvio.

\subsection{Monitor e \emph{baseline} com verdade conhecida}
\label{sec:verdade-conhecida}
% [REV-2026-09-13][C7-14][P1][SIMULACAO]
\revisaotese{C7-14}{P1}{SIMULACAO}
{Tornar o gerador e a calibração de falsos sinais comparáveis}
{Especificar equação do gerador, escala marginal versus escala da
inovação AR(1), covariáveis, onset, sinal, lateralidade e tratamento de
lacunas. "5\%" é por dia, canal ou probabilidade de algum sinal em 90 dias?
Usar o mesmo orçamento por horizonte na comparação e réplicas
independentes para afinar limiares e os verificar. Se combinar canais,
calibrar também o sistema. Dar incerteza Monte Carlo; um canal acima
do nominal não se declara "acaso" sem o critério global pré-definido.}%
% [/REV-2026-09-13]

A pergunta que o calendário operacional tornou não estimável no
Capítulo~\ref{cha:diagnostico} --- que desvio o sistema deteta, e o que
importa mais para o detetar --- só pode responder-se com verdade conhecida.
Construiu-se, para cada vetor, um gerador com a relação de 2020 (mínimos
quadrados com as covariáveis de referência), ruído autorregressivo de
primeira ordem com $\rho = 0{,}6$ e dispersão igual à dos resíduos de 2020,
366 dias de treino com a máscara real de dias admissíveis de 2020 e uma
aplicação de noventa dias; os limiares de cada monitor foram calibrados sob
a hipótese nula a uma taxa de falso sinal de $5\%$ e verificados (quarenta
canais entre $0{,}020$ e $0{,}072$, um acima do nominal pelo limite de
Wilson~\cite{Wilson1927}, compatível com o acaso). Sobre este gerador inseriram-se
alternativas conhecidas --- degrau e rampa de uma escala, alargamento da
dispersão, recuperação, resposta à carga --- em quatrocentas réplicas por
célula, e mediu-se a taxa de deteção de cada combinação de \emph{baseline}
(mínimos quadrados ou modelo aditivo~\cite{HastieTibshirani1986,Wood2017}) e
monitor (resíduo bruto, filtro de
nível, soma acumulada, e canais de carga e de escala). \textbf{Nada nesta
secção é um facto sobre a refinaria}: é o comportamento de métodos sobre um
gerador que escrevemos, e deve ser lido como tal.

\begin{table}[htbp]
  \centering
  \caption[Contrastes co-primários com verdade conhecida]{Diferença de taxa
  de deteção de um degrau de $1\sigma$, emparelhada por réplica, com intervalo
  percentil de 2\,000 reamostragens. Esquerda: trocar a \emph{baseline}
  (mínimos quadrados $\to$ modelo aditivo) com o monitor de nível fixo.
  Direita: trocar o monitor (filtro de nível $\to$ soma acumulada) com a
  \emph{baseline} fixa. Gerador com $\rho = 0{,}6$.}
  \label{tab:coprimarias}
  \begin{tabular}{lcccc}
    \toprule
    Vetor & Trocar a \emph{baseline} & IC 95\% & Trocar o monitor & IC 95\% \\
    \midrule
    Fuel gás     & $-0{,}028$ & $[-0{,}052;\,-0{,}005]$ & $+0{,}323$ & $[0{,}273;\,0{,}370]$ \\
    Vapor 24 bar & $-0{,}013$ & $[-0{,}040;\,0{,}015]$  & $+0{,}438$ & $[0{,}388;\,0{,}487]$ \\
    Vapor 10 bar & $-0{,}020$ & $[-0{,}045;\,0{,}003]$  & $+0{,}490$ & $[0{,}440;\,0{,}540]$ \\
    Vapor 3 bar  & $-0{,}028$ & $[-0{,}060;\,0{,}003]$  & $+0{,}390$ & $[0{,}343;\,0{,}438]$ \\
    \bottomrule
  \end{tabular}
\end{table}

% [REV-2026-09-13][C7-15][P2][SIMULACAO]
\revisaotese{C7-15}{P2}{SIMULACAO}
{O gerador linear favorece a comparação pontual adotada}
{Como o gerador primário usa relação linear, pequeno ganho do GAM não
demonstra que a escolha da baseline importe pouco em geral.
Dar diferenças absolutas (32--49 pontos percentuais), IC e estatuto
por contraste; oito testes co-primários sem correção não são evidência
simultânea de todos os efeitos. Reportar gerador não linear como
sensibilidade, sem o promover após observar resultados. A secção já
limita a interpretação: aplicar a mesma cautela ao título e às conclusões.}%
% [/REV-2026-09-13]
\paragraph{O monitor pesa; a \emph{baseline} não --- nas oito células
declaradas.} As oito comparações declaradas como primárias antes da
execução (Tabela~\ref{tab:coprimarias}) dizem, nos quatro vetores, a mesma
coisa: trocar o monitor de nível por uma soma acumulada compra entre um terço
e metade da taxa de deteção de um degrau de uma escala, com os quatro
intervalos longe de zero; trocar a \emph{baseline} linear por um modelo
aditivo custa um a três pontos, e só no fuel gás o intervalo exclui zero ---
a favor da linear. Os quatrocentos e quarenta contrastes secundários,
publicados sem nunca serem promovidos, dizem onde isto \emph{não} vale: num
alargamento puro de dispersão a soma acumulada \emph{perde} ($-0{,}28$ em
média), porque acumula desvios de média e esses cancelam-se; em rampas e
recuperações compra $0{,}29$ e $0{,}20$; e num banco com interação
carga$\times$densidade a \emph{baseline} aditiva ganha. A escolha do monitor
é um compromisso, não uma melhoria --- e nenhuma recomendação vale sem a
natureza do desvio e a autocorrelação medidas ao lado.

% [REV-2026-09-13][F20][P2][FIGURA]
\revisaotese{F20}{P2}{FIGURA}
{Curvas de deteção e atraso sob o mesmo falso sinal}
{Usar painéis para rho=0,6/0,9/0,95, amplitude do desvio no eixo x e
P(deteção até 90 dias) no y, com IC Monte Carlo e um traço por monitor.
Marcar censura e limiar de 50\%; separar degrau de rampa.
Acrescentar distribuição do atraso com não deteções censuradas, se
calculada; reportar só atraso entre detetados favorece métodos pouco
sensíveis. Identificar gerador simulado, não previsão para a refinaria.}%
% [/REV-2026-09-13]
\paragraph{A dependência decide quanto se deteta.} O gerador primário tem
$\rho = 0{,}6$, que é a autocorrelação residual observada no fuel gás; os
três vapores estão entre $0{,}8$ e $0{,}9$, e a subir
(Secção~\ref{sec:diagnosticos-desc}). O lote de sensibilidade responde ao que
isso muda. Num degrau de uma escala, a soma acumulada deteta $0{,}96$ das
réplicas a $\rho = 0{,}6$, $0{,}59$ a $\rho = 0{,}9$ e $0{,}39$ a
$\rho = 0{,}95$ --- neste último caso igual ao resíduo bruto ($0{,}35$): a
vantagem do monitor desaparece. Numa rampa, a taxa cai de $0{,}64$ para
$0{,}27$ e $0{,}17$. A leitura para os vapores é, por isso, mais modesta do
que a das oito células primárias: na dependência que eles exibem, um degrau
de uma escala é detetado dentro do horizonte em cerca de metade das
réplicas pelo melhor monitor, e uma rampa em um quarto. O monitor continua
a pesar; a detetabilidade absoluta é pobre, e o resultado do fuel gás não
transfere.

\paragraph{O desvio mínimo detetável, no gerador.} Sobre a curva de deteção
por amplitude de degrau, o desvio a que a soma acumulada deteta metade das
vezes situa-se entre $0{,}51\sigma$ (vapor de 3~bar) e $0{,}63\sigma$ (vapor
de 10~bar), contra $0{,}87$ a $1{,}10\sigma$ do filtro de nível e $0{,}95$ a
$1{,}15\sigma$ do resíduo bruto. Em percentagem do consumo médio dos dias
retidos de 2020, o desvio mínimo detetável da soma acumulada é $4{,}5\%$ no
fuel gás, $9{,}4\%$ no vapor de 3~bar e $14{,}7\%$ no de 24~bar; o vapor de
10~bar é um fluxo líquido com sinal e não admite percentagem. Trinta e duas
das oitenta curvas estão censuradas à direita --- todas nos canais de carga e
de escala, que não existem para ver um degrau de nível --- e são publicadas
como limite aberto, nunca como número. Estes valores são propriedades do
gerador a $\rho = 0{,}6$; na refinaria, e nos vapores, o parágrafo anterior
diz o que se pode esperar.

% [REV-2026-09-13][C7-16][P2][SIMULACAO]
\revisaotese{C7-16}{P2}{SIMULACAO}
{Exemplos simulados não são teoremas universais}
{Dar médias/tendências das séries independentes e distribuição de R\textsuperscript{2}
em várias réplicas; R\textsuperscript{2} acumulado alto nesta realização não torna toda
relação acumulada espúria. Branquear pode atenuar um degrau, mas não o
remove necessariamente (em AR(1) fixo, a resposta sustentada depende
de 1-rho); monitorizar inovações não é o mesmo que threshold diário.
Explicar 0,095 como taxa nula face ao alvo 0,05 referido acima e confirmar
que canais conjuntos/marginais usam lateralidade e falso sinal comparáveis.}%
% [/REV-2026-09-13]
\paragraph{Três sinais falsos, demonstrados.} O mesmo aparelho serviu para
demonstrar, com verdade conhecida, três leituras sedutoras que a prática de
monitorização produz. Primeira: duas séries \emph{independentes} têm $R^2$
diário de $0{,}0025$, mensal de $0{,}18$ e acumulado de $0{,}9999$ --- a
correlação de séries acumuladas é inteiramente fabricada pela
acumulação~\cite{Yule1926,GrangerNewbold1974}, e
um gráfico de consumo acumulado contra carga acumulada não é evidência de
relação nenhuma. Segunda: branquear a série antes de monitorizar --- usar as
inovações de um modelo autorregressivo como estatística --- remove
exatamente a mudança de nível que se quer ver~\cite{AlwanRoberts1988,Wardell1994}: num degrau de uma escala, a
inovação autorregressiva sinaliza à sua taxa nula ($0{,}095$), enquanto uma
média móvel de trinta dias deteta $0{,}895$. Terceira: a leitura de que a
anticorrelação entre os vapores de 24 e 10~bar
(Secção~\ref{sec:diagnosticos-desc}) ``impede a coocorrência'' de alarmes ---
que constou de uma versão anterior desta exploração e foi retirada ---
assentava em pontuar o vapor de 10~bar como unilateral; na lateralidade
declarada, com os quatro vetores desviados, um canal conjunto sinaliza em
$0{,}068$ das réplicas contra $0{,}570$ da melhor margem individual, e com um
só vetor desviado em $0{,}022$. Uma contagem conjunta exige mais do que a
melhor margem e por isso deteta menos; não é a correlação que o impede.

\paragraph{O que a exploração com referência congelada ainda não fecha.}
Esta exploração foi submetida a duas revisões independentes. A primeira
aprovou a integridade da cadeia de resultados e reprovou a conformidade
científica de sete pontos, todos corrigidos e republicados. A segunda,
sobre a versão republicada, identificou que várias das suas partes --- a
comparação entre esquemas de referência, a ordenação de janelas de
reajuste, a inserção de desvios sobre estados publicados, a agregação a
frequências semanal e mensal e o suporte da eletricidade --- ainda não
estimam exatamente a comparação declarada (comparam populações diferentes,
comprimem o calendário, ou incluem extrapolações no resultado primário). As
correções estão em curso à data desta versão; este capítulo cita, com
números, apenas os lotes que essa revisão classificou como aceitáveis, e
onde uma parte em revisão entra, entra qualitativamente e assinalada.

\section{Implicações para a plataforma}
\label{sec:implicacoes-plataforma}
% [REV-2026-09-13][F21][P3][FIGURA]
\revisaotese{F21}{P3}{FIGURA}
{Protótipo da ficha e diagrama de estados}
{Após os requisitos, mostrar um wireframe identificado como "proposta",
com referência/data, domínio, desvio, cobertura/IC, estado e ação.
Uma segunda faixa representa operação -\textgreater{} investigar -\textgreater{} suspender/
recalibrar -\textgreater{} verificar, com gatilhos e responsável; separar eventos
(alarme) de estados persistentes. A versão anterior da referência
permanece consultável. Usar poucos campos legíveis na ficha principal;
os 76 campos e o registo detalhado ficam no Apêndice E.}%
% [/REV-2026-09-13]

O Capítulo~\ref{cha:plataforma} entregou um instrumento que torna o
diagnóstico possível; este capítulo devolve-lhe uma especificação. O que se
segue é o bloco de saúde das \emph{baselines} --- o que o \emph{dashboard}
tem de mostrar para que uma \emph{baseline} seja usada com conhecimento do
seu estado --- formulado como requisitos, cada um ancorado num resultado
desta dissertação.

\paragraph{Duas pistas, não uma.} O método de produção usa um só modelo para
normalizar e para alarmar. A plataforma deve manter duas pistas visíveis
lado a lado: o modelo anual em vigor, para reporte normalizado e comparação
entre períodos, como a ISO~50006 pede; e uma referência congelada, com o
desvio padronizado face a ela e um monitor sequencial sobre esse desvio,
para deteção. A Secção~\ref{sec:referencia-fixa} mostrou que só a segunda
mede o desvio; a Secção~\ref{sec:verdade-conhecida} mostrou que o monitor
sobre ela é o que mais pesa na deteção. A data e a razão de cada mudança de
referência são registadas.

\paragraph{Domínio visível, alarme suprimido fora dele.} Cada
\emph{baseline} publica o intervalo de carga e as condições em que foi
estimada; um dia fora desse domínio é mostrado como ``fora de domínio'', não
como anomalia, e o seu alarme é suprimido. O limiar de carga oficial já
define um domínio implícito (Secção~\ref{sec:perfil-dados}); a plataforma
torna-o explícito.

% [REV-2026-09-13][C7-13][P1][INTERPRETACAO]
\revisaotese{C7-13}{P1}{INTERPRETACAO}
{Cobertura e taxa de excursões são complementares}
{Para a mesma banda, população e regra, cobertura = 1 - taxa de dias
fora da banda. Trocar uma pela outra não elimina confusão entre
descalibração e desvio; ambas observam o mesmo evento. A incerteza e o
contexto operacional é que acrescentam informação. Especificar taxa
por dia versus taxa por episódio, superior/inferior, e usar investigação
externa para discriminar causas. Janelas móveis sobrepostas também
não constituem testes independentes repetidos.}%
% [/REV-2026-09-13]
\paragraph{Calibração monitorizada, não assumida.} A cobertura efetiva da
banda de alarme nos últimos noventa e cento e oitenta dias, com intervalo de
confiança que respeite a dependência, é um indicador do bloco de saúde ---
não a taxa de alarmes, que confunde descalibração com desvio. Uma cobertura
que exclui a nominal pelo intervalo desencadeia a política de investigação,
não uma recalibração automática.

\paragraph{Ficha por \emph{baseline}.} Os setenta e seis campos da
Secção~\ref{sec:demo-fichas} --- domínio, janela, exclusões com códigos de
razão, estimador, método de banda e desfecho, resultados do critério de
aceitação, proveniência --- são a ficha que acompanha cada \emph{baseline}
na plataforma, gerada dos resultados e nunca editada à mão. Uma
\emph{baseline} sem ficha não está em produção.

\paragraph{Políticas separadas e registadas.} Alarme, investigação, suspensão,
recalibração e verificação pós-recalibração são cinco estados com transições
declaradas; a recalibração exige sinal do monitor mais revisão documentada,
com a opção ``causa desconhecida, recalibrar com sinalização''; a
certificação do modelo recalibrado usa só dados posteriores. O registo de
versões, aprovações e revisões é parte do bloco, não um documento à parte.

\paragraph{O que a plataforma não pode fazer sem informação que não tem.}
Nenhum dos requisitos acima classifica um desvio como ineficiência, nem um
alarme como falso. Isso exige o calendário operacional, rótulos de eventos e
margens materiais elicitadas com a operação --- as três coisas que esta
dissertação não teve. A plataforma pode e deve registar o desvio e o alarme
com a sua proveniência; a classificação é uma decisão humana até que essa
informação exista, e o desenho deve deixar espaço para a receber.

\section{Limitações}
\label{sec:limitacoes}
% [REV-2026-09-13][C7-17][P2][LIMITACOES]
\revisaotese{C7-17}{P2}{LIMITACOES}
{Corrigir contradições internas e preservar o alcance real}
{Distinguir previsão temporal fora do treino de validação independente
de seleção: os pares são forward, mas o estudo já usou esses anos para
formular decisões. Evitar "não há afirmação fora de amostra que
sobreviva". Conferir 3/16 na sensibilidade da Sec. 7.4 contra 3/32 aqui.
O ensaio de cobertura do Cap. 4 usa bloco 60; não o atribuir ao bloco
base sete. Trocar "resposta completa" e "máximo defensável" por opções
de investigação: os resultados não demonstram necessidade nem suficiência
de modelos com estado de degradação.}%
% [/REV-2026-09-13]

\paragraph{Do estudo.} Uma unidade, seis anos completos, um único
\emph{snapshot} de dados já explorados: não há afirmação fora de amostra que
sobreviva, e o desenho não a tenta --- a demonstração das linhas
orientadoras é retrospetiva e encadeada, nunca \emph{holdout}, porque 2025 e
o parcial de 2026 já tinham informado o diagnóstico. As conclusões são
condicionais ao modelo linear, à seleção oficial e às janelas anuais que o
método de produção herdou. O calendário operacional, os rótulos de eventos e
as margens materiais não estavam disponíveis: nenhum desvio é ineficiência,
nenhuma excursão é falso alarme, nenhuma cobertura é qualidade de deteção,
e a detetabilidade só foi estudada em simulação. Há erro de medição na
carga, possível endogeneidade (a operação reage aos alarmes), agregação
diária que esconde dinâmica intra-diária, um mecanismo de falta não
identificável sem fonte operacional, e nenhum padrão de referência
independente do modelo Power BI original para as observações selecionadas.
A energia elétrica é uma alocação mensal e fica fora de qualquer inferência
diária.

\paragraph{Da inferência.} Todos os veredictos assentam em cinco pares
anuais e em regras binárias sobre eles; a potência dessas regras sob a
dependência observada não foi calculada, e o zero em trinta e dois passa a
três em trinta e dois com uma regra de três em cinco. O intervalo de
confiança por \emph{bootstrap} de blocos usa um bloco base de sete dias
sobre resíduos com autocorrelação até $0{,}90$; a sua cobertura nominal não
está demonstrada --- a validação por simulação a $\rho = 0{,}90$ deu
$0{,}8925$ (intervalo de Wilson $0{,}86$ a $0{,}92$), o que não é evidência
de cobertura inferior a $0{,}90$ nem prova de cobertura de $0{,}95$. Um
contraste que exclui zero por pouco está dentro do erro de calibração do
próprio instrumento, e a sensibilidade ao comprimento de bloco reportada
nas Secções~\ref{sec:estabilidade} e~\ref{sec:calibracao} é o sintoma
visível disso. A correção de multiplicidade aplica-se só à família dos
testes de estabilidade; todos os outros resultados são por contraste, e a
sua contagem não é evidência conjunta. ``Não rejeitado'' nunca é
``equivalente''~\cite{Lakens2017}: não há margens práticas elicitadas.

\paragraph{Da simulação.} O gerador primário tem $\rho = 0{,}6$, o regime do
fuel gás, não o dos vapores; depende da relação de referência, das
máscaras, do ruído e do orçamento de réplicas escolhidos; o modelo aditivo
corre com hiperparâmetros fixos; o desvio mínimo detetável é do gerador, não
da unidade; e um número obtido em quatrocentas réplicas de um gerador que
escrevemos não é um facto sobre a refinaria, por muitas réplicas que tenha.

\paragraph{Das linhas orientadoras.} A classe linear robustificada, com
bandas honestas e critério de aceitação, leva o modelo estático ao máximo
defensável --- e não resolve a não-estacionariedade estrutural que o
diagnóstico documenta. O critério de aceitação é um procedimento com
critérios declarados, demonstrado com $n = 5$; o que ele decide depende do
limiar, e o que não decide é a impossibilidade da classe. A resposta
completa --- modelos com estado de degradação, referência congelada com
monitor sequencial em produção, validação com eventos reais --- é a
investigação seguinte (Capítulo~\ref{cha:conclusoes}), para a qual ficam
prontos a plataforma, o diagnóstico, o critério de aceitação e a evidência
auditável que os sustenta.
